\documentclass[12pt,a4paper]{article}

% ---------------- PACKAGES ----------------
\usepackage[utf8]{inputenc}
\usepackage[margin=1in]{geometry}
\usepackage{graphicx}
\usepackage{amsmath}
\usepackage{amssymb}
\usepackage{times}
\usepackage{setspace}
\usepackage{ragged2e}
\usepackage{enumitem}
\usepackage[most]{tcolorbox}
\usepackage{listings} % কোড বক্স মার্জিনের ভেতরে রাখার জন্য নতুন প্যাকেজ

% ---------------- SETTINGS ----------------
\setstretch{1.5}
\justifying

% কোড বক্স কনফিগারেশন (যা লম্বা কোডকে অটোমেটিক ভেঙে বক্সের ভেতরে রাখবে)
\lstnewenvironment{codebox}{
\lstset{
    basicstyle=\small\ttfamily,
    border=single,
    breaklines=true,         % লম্বা লাইন অটোমেটিক ভেঙে নিচে নামবে
    breakatwhitespace=false, % যেকোনো জায়গায় লাইন ভাঙতে পারবে
    frame=single,
    frameround=ffff,
    framesep=5mm,
    showstringspaces=false,
    columns=fullflexible
}
}{}

% সূচিপত্রের কাস্টম কমান্ড (পেজ নাম্বার ছাড়া পরিচ্ছন্ন ডিজাইন)
\newcommand{\tocline}[1]{\vspace{0.25cm}\noindent\textbf{\large #1}\\[0.1cm]}
\newcommand{\tocsub}[1]{\noindent\hspace*{6mm}\textbullet~~#1\\[0.05cm]}

\begin{document}

% ---------------- TITLE PAGE (কাভার পেজ) ----------------
\begin{titlepage}
\centering
{\Large\bfseries Begum Rokeya University, Rangpur}\\[0.4cm]
{\large\bfseries Department of Statistics}\\[2.5cm]
\vfill
{\large\bfseries Course Code: 4101}\\[0.3cm]
{\LARGE\bfseries Categorical Data Analysis}\\[3cm]
\vfill

\begin{minipage}[t]{0.48\textwidth}
\raggedright
{\large\bfseries SUBMITTED BY}\\[0.2cm]
\textbf{Name:} Md. Sifat Hasan Sabbir\\
\textbf{ID:} 12110045\\
\textbf{Reg No:} 000015908\\
\textbf{Session:} 2021-2022\\
\textbf{Year/Sem:} 4th Year, 1st Sem
\end{minipage}
\hfill
\begin{minipage}[t]{0.48\textwidth}
\raggedleft
{\large\bfseries SUBMITTED TO}\\[0.2cm]
\textbf{Name:} Dr. Md. Siddikur Rahman\\
Associate Professor\\
Department of Statistics\\
Begum Rokeya University, Rangpur
\end{minipage}

\vfill
\hrule
\vspace{0.3cm}
{\large\bfseries Submission Date: 19/05/2026}
\end{titlepage}

% ---------------- TABLE OF CONTENTS (সূচিপত্র) ----------------
\newpage
\begin{center}
    \vspace*{0.2cm}
    {\Large\bfseries সূচিপত্র (Table of Contents)} \\[0.2cm]
    \hrule height 1.5pt
\end{center}
\vspace{0.5cm}

\tocline{Problem-1: Categorical Data Analysis \& Lung Disease}
\tocsub{ Basic EDA Questions}
\tocsub{ Association & Crosstab}
\tocsub{ Does Pollution level affect Lung Disease?}
\tocsub{ Statistical Testing}
\tocsub{ Correlation & Interpretation}
\tocsub{ Logistic Regression}
\tocsub{ Model Evaluation}
\tocsub{Final Analysis and Public Health Implications}


\tocline{Problem-2: One-Way ANOVA \& Weight Loss Analysis}
\tocsub{Assumption Checking (Normality, Homogeneity) \& ANOVA Test}

\tocline{Problem-3: Contingency Table Analysis (Anemia vs Wealth Index)}
\tocsub{NDHS Data Table, Chi-Square Test \& Heatmap Visualization}

\tocline{Problem-4: Drug Effectiveness Analysis}
\tocsub{Global Association \& Post-hoc Pairwise Fisher's Exact Tests}

\tocline{Problem-5: Logistic Regression Analysis (Graduate Admission)}
\tocsub{Model Estimation (GRE, GPA, Rank) \& Odds Ratio (OR) Analysis}

\tocline{Problem-6: Poisson Regression Analysis for Count Data}
\tocsub{GLM Model Fitting, Equations \& Incident Rate Ratios (IRR)}

\tocline{Problem-7: Negative Binomial Regression for Absenteeism}
\tocsub{Overdispersion Diagnosis \& Negative Binomial Modeling}

\tocline{Problem-8: Zero-Inflation Analysis (Fish Caught Study - Part I)}
\tocsub{Excess Zeros Frequency Analysis, ZIP Model Coefficients \& IRR}

\tocline{Problem-9: Zero-Inflated Poisson Regression (Part II)}
\tocsub{Exploratory Data Analysis, Logit vs. Poisson Component \& IRR}

\tocline{Problem-10: Negative Binomial for Hospital Length of Stay}
\tocsub{Zero-Truncated Context (ZTNB Model) \& IRR Interpretation}

\newpage



% ---------------- ASSIGNMENT ----------------
\begin{center}
\Large \textbf{Problem-1}
\end{center}
\textbf{A: Basic EDA Questions}
% ==================================================
% ---------------- i ----------------
% ==================================================

\textbf{i)}

\textbf{Code:}
\begin{codebox}
import os
import pandas as pd
import matplotlib.pyplot as plt
import seaborn as sns

file_path = r"C:\Users\hp\OneDrive\Desktop
\Reality Stone-4.1\Siddik sir-4104\lung_disease.xlsx"
save_dir = r"C:\Users\hp\OneDrive\Desktop
\Reality Stone-4.1\Siddik sir-4104"
df = pd.read_excel(file_path)
sns.set_theme(style="whitegrid")

age_summary = df['Age'].describe()
age_skewness = df['Age'].skew()

print("--- Age Distribution Summary ---")
print(age_summary)
print(f"Skewness of Age: {age_skewness:.4f}")

plt.figure(figsize=(8,5))
sns.histplot(df['Age'], kde=True, bins=20)
plt.title('Distribution of Age')
plt.xlabel('Age')
plt.ylabel('Count')
plt.savefig(os.path.join(save_dir, 
'age_distribution.png'), dpi=300, bbox_inches='tight')
plt.show()
\end{codebox}

\textbf{Output:}
\begin{codebox}
--- Age Distribution Summary ---
count    500.000000
mean      43.904000
std       14.604841
min       20.000000
25%       30.750000
50%       45.000000
75%       56.000000
max       69.000000
Skewness of Age: 0.0121
\end{codebox}

\textbf{Graph:}
\begin{center}
\includegraphics[width=0.7\textwidth]
{age_distribution.png}
\end{center}

\textbf{Interpretation:}
The Age variable shows an approximately symmetric distribution with mean 43.90 and standard deviation 14.60.

% ==================================================
% ---------------- ii ----------------
% ==================================================
\textbf{ii)}
\textbf{Code:}
\begin{codebox}
smoking_counts = df['Smoking'].value_counts()
smoking_proportions = df['Smoking'].
value_counts(normalize=True)

print("--- Smoking Counts and Proportions ---")
for idx in smoking_counts.index:
    print(f"{idx}: {smoking_counts[idx]} individuals ({smoking_proportions[idx]*100:.2f}%)")

plt.figure(figsize=(6,6))
plt.pie(smoking_counts, labels=smoking_counts.index,
autopct='%1.1f%%',
        startangle=90, colors=['#ff9999','#66b3ff'])
plt.title('Proportion of Smokers vs Non-Smokers')
plt.savefig(os.path.join(save_dir, 
'smoking_proportion.png'), dpi=300,
bbox_inches='tight')
plt.show()
\end{codebox}

\textbf{Output:}
\begin{codebox}
--- Smoking Counts and Proportions ---
No: 293 individuals (58.60%)
Yes: 207 individuals (41.40%)
\end{codebox}

\textbf{Graph:}
\begin{center}
\includegraphics[width=0.6\textwidth]
{smoking_proportion.png}
\end{center}

\textbf{Interpretation:}
The dataset shows 58.60 percent non-smokers and 41.40 percent smokers.

% ==================================================
% ---------------- iii ----------------
% ==================================================
\textbf{iii)}
\textbf{Code:}
\begin{codebox}
income_counts = df['Income'].value_counts()

print("--- Income Group Frequency ---")
print(income_counts)
print(f"\nHighest frequency income group is:
{income_counts.idxmax()}")

plt.figure(figsize=(7,5))
sns.countplot(x='Income', data=df,
order=['Low', 'Medium', 'High'], palette='viridis')
plt.title('Frequency of Income Groups', fontsize=14)
plt.xlabel('Income Group', fontsize=12)
plt.ylabel('Count', fontsize=12)
plt.savefig(os.path.join(save_dir,
'income_frequency.png'), dpi=300, bbox_inches='tight')
plt.show()
\end{codebox}

\textbf{Output:}
\begin{codebox}
--- Income Group Frequency ---
Income
Low       204
Medium    197
High       99
Name: count, dtype: int64
Highest frequency income group is: Low
\end{codebox}

\textbf{Graph:}
\begin{center}
\includegraphics[width=0.7\textwidth]{income_frequency.png}
\end{center}

\textbf{Interpretation:}
The income distribution shows Low income group is highest (204), followed by Medium (197), and High (99).

% ==================================================
% ---------------- iv ----------------
% ==================================================
\textbf{iv)}
\textbf{Code:}
\begin{codebox}
pollution_proportions = df['Pollution'].
value_counts(normalize=True) * 100

print("--- Pollution Exposure Percentage ---")
print(pollution_proportions.map('{:.2f}%'.format))

plt.figure(figsize=(6,5))
sns.countplot(x='Pollution', data=df, order=['High', 'Low'], palette='magma')
plt.title('Pollution Exposure Level', fontsize=14)
plt.xlabel('Pollution Level', fontsize=12)
plt.ylabel('Count', fontsize=12)
plt.savefig(os.path.join(save_dir, 
'pollution_exposure.png'), dpi=300, bbox_inches='tight')
plt.show()
\end{codebox}

\textbf{Output:}
\begin{codebox}
--- Pollution Exposure Percentage ---
Pollution
High    70.40%
Low     29.60%
Name: proportion, dtype: object
\end{codebox}

\textbf{Graph:}
\begin{center}
\includegraphics[width=0.7\textwidth]{pollution_exposure.png}
\end{center}

\textbf{Interpretation:}
The majority (70.40\%) of individuals are exposed to high pollution levels.

% ==================================================
% ---------------- v ----------------
% ==================================================
\textbf{v)}
\textbf{Code:}
\begin{codebox}
disease_counts = df['LungDisease'].value_counts()
disease_prevalence = df['LungDisease']
.value_counts(normalize=True) * 100

print("--- Lung Disease Prevalence ---")
print(f"Yes: {disease_counts['Yes']} individuals ({disease_prevalence['Yes']:.2f}%)")
print(f"No: {disease_counts['No']} individuals ({disease_prevalence['No']:.2f}%)")

plt.figure(figsize=(6,6))
plt.pie(disease_counts, labels=disease_counts.index, 
autopct='%1.1f%%',
        startangle=140, colors=['#ffcc99','#99ff99'], wedgeprops=dict(width=0.4, edgecolor='w'))
plt.title('Overall Prevalence of Lung Disease')
plt.savefig(os.path.join(save_dir, 
'lung_disease_prevalence.png'), dpi=300, bbox_inches='tight')
plt.show()
\end{codebox}

\textbf{Output:}
\begin{codebox}
--- Lung Disease Prevalence ---
Yes: 264 individuals (52.80%)
No: 236 individuals (47.20%)
\end{codebox}

\textbf{Graph:}
\begin{center}
\includegraphics[width=0.6\textwidth]
{lung_disease_prevalence.png}
\end{center}

\textbf{Interpretation:}
The prevalence of lung disease is 52.80\% Yes and 47.20\% No.

% ==================================================
% ---------------- B ----------------
% ==================================================
\textbf{B. Association \& Crosstab}
\textbf{i)}
\textbf{Code:}
\begin{codebox}
import os
import pandas as pd
import matplotlib.pyplot as plt
import seaborn as sns
import scipy.stats as stats

file_path = r"C:\Users\hp\OneDrive\Desktop
\Reality Stone-4.1\Siddik sir-4104\lung_disease.xlsx"
save_dir = r"C:\Users\hp\OneDrive\Desktop
\Reality Stone-4.1\Siddik sir-4104"
df = pd.read_excel(file_path)

contingency_table = pd.crosstab(df['Smoking'],
df['LungDisease'])
chi2_stat, p_val, dof, expected = 
stats.chi2_contingency(contingency_table)

print("--- Chi-Square Test Results ---")
print(f"Chi-square Statistic : {chi2_stat:.4f}")
print(f"p-value : {p_val:.4e}")
print(f"Degrees of Freedom : {dof}")

plt.figure(figsize=(7,5))
sns.countplot(x='Smoking', hue='LungDisease',
data=df, palette='Set2')
plt.title('Association between Smoking and
Lung Disease', fontsize=14)
plt.xlabel('Smoking Status', fontsize=12)
plt.ylabel('Count', fontsize=12)
plt.legend(title='Lung Disease')
plt.savefig(os.path.join(save_dir, 
'smoking_lung_disease_association.png'), dpi=300, bbox_inches='tight')
plt.show()
\end{codebox}

\textbf{Output:}
\begin{codebox}
--- Chi-Square Test Results ---
Chi-square Statistic : 67.5943
p-value : 2.0085e-16
Degrees of Freedom : 1
\end{codebox}

\textbf{Graph:}
\begin{center}
\includegraphics[width=0.7\textwidth]
{smoking_lung_disease_association.png}
\end{center}

\textbf{Interpretation:}
The Chi-square test shows a statistically significant association between smoking status and lung disease (p-value $<$ 0.05).
This indicates that smoking and lung disease are not independent.
The high Chi-square value (67.5943) suggests a strong relationship between the two variables, meaning smokers are more likely to have lung disease compared to non-smokers.

% ---------------- ii ----------------
\textbf{ii)}
\textbf{Code:}
\begin{codebox}
import os
import pandas as pd
import matplotlib.pyplot as plt
import seaborn as sns

file_path = r"C:\Users\hp\OneDrive\Desktop
\Reality Stone-4.1\Siddik sir-4104\lung_disease.xlsx"
save_dir = r"C:\Users\hp\OneDrive\Desktop
\Reality Stone-4.1\Siddik sir-4104"
df = pd.read_excel(file_path)

crosstab_counts = pd.crosstab(df['Smoking'],
df['LungDisease'])
crosstab_pct = pd.crosstab(df['Smoking'],
df['LungDisease'], normalize='index') * 100

print("--- Contingency Table (Counts) ---")
print(crosstab_counts)
print("\n--- Contingency Table (Percentages) ---")
print(crosstab_pct.map('{:.2f}%'.format))

plt.figure(figsize=(6,4))
sns.heatmap(crosstab_counts, annot=True, fmt='d',
cmap='Blues', cbar=False)
plt.title('Contingency Table Heatmap (Counts)',
fontsize=14)
plt.xlabel('Lung Disease', fontsize=12)
plt.ylabel('Smoking', fontsize=12)
plt.savefig(os.path.join(save_dir, 
'contingency_table_heatmap.png'), dpi=300, bbox_inches='tight')
plt.show()
\end{codebox}

\textbf{Output:}
\begin{codebox}
--- Contingency Table (Counts) ---
LungDisease    No   Yes
Smoking
No            184   109
Yes            52   155
--- Contingency Table (Percentages) ---
LungDisease      No     Yes
Smoking
No           62.80%  37.20%
Yes          25.12%  74.88%
\end{codebox}

\textbf{Graph:}
\begin{center}
\includegraphics[width=0.65\textwidth]{contingency_table_heatmap.png}
\end{center}

\textbf{Interpretation:}
The contingency table shows a strong relationship between smoking and lung disease.
Among non-smokers, 37.20\% have lung disease, whereas among smokers, 74.88\% have lung disease.
The heatmap visually confirms this pattern, showing higher concentration of lung disease among smokers.
This indicates a strong positive association between smoking and lung disease.

% ==================================================
% ---------------- C ----------------
% ==================================================
\textbf{C. Income Analysis and Factor Strength}
\textbf{i) Does Income level affect Lung Disease?}
\textbf{Code:}
\begin{codebox}
import os
import pandas as pd
import matplotlib.pyplot as plt
import scipy.stats as stats

file_path = r"C:\Users\hp\OneDrive\Desktop
\Reality Stone-4.1\Siddik sir-4104\lung_disease.xlsx"
save_dir = r"C:\Users\hp\OneDrive\Desktop
\Reality Stone-4.1\Siddik sir-4104"
df = pd.read_excel(file_path)

income_xtab = pd.crosstab(df['Income'], df['LungDisease'])
income_pct = pd.crosstab(df['Income'], df['LungDisease'], normalize='index') * 100

print("--- Income vs Lung Disease (Counts) ---")
print(income_xtab)
print("\n--- Income vs Lung Disease
(Prevalence Percentages) ---")
print(income_pct.map('{:.2f}%'.format))

chi2_inc, p_val_inc, _, _ = stats.chi2_contingency(income_xtab)
print(f"\nChi-Square Statistic: {chi2_inc:.4f}")
print(f"Chi-Square p-value : {p_val_inc:.4f}")

income_pct.plot(kind='bar', stacked=True,
color=['#4caf50', '#f44336'], figsize=(8,5))
plt.title('Lung Disease Prevalence across
Income Groups', fontsize=14, fontweight='bold')
plt.xlabel('Income Group', fontsize=12)
plt.ylabel('Percentage (%)', fontsize=12)
plt.legend(title='Lung Disease',
labels=['No', 'Yes'], loc='upper left')
plt.xticks(rotation=0)
plt.savefig(os.path.join(save_dir, 
'income_vs_lung_disease.png'),
dpi=300, bbox_inches='tight')
plt.show()
\end{codebox}

\textbf{Output:}
\begin{codebox}
--- Income vs Lung Disease (Counts) ---
LungDisease   No  Yes
Income
High          51   48
Low           90  114
Medium        95  102
--- Income vs Lung Disease (Prevalence Percentages) ---
LungDisease      No     Yes
Income
High         51.52%  48.48%
Low          44.12%  55.88%
Medium       48.22%  51.78%
Chi-Square Statistic: 1.6002
Chi-Square p-value : 0.4493
\end{codebox}

\textbf{Interpretation:}
The Chi-square test shows that there is no statistically significant association between income level and lung disease (p-value = 0.4493 $>$ 0.05).
Although slight differences exist in disease prevalence across income groups, these differences are not strong enough to conclude a meaningful relationship.
Therefore, income does not appear to be a significant predictor of lung disease in this dataset.

% ---------------- ii ----------------
\textbf{ii) Which factor appears to have the strongest relationship with lung disease?}
\textbf{Code:}
\begin{codebox}
import os
import pandas as pd
import numpy as np
import matplotlib.pyplot as plt
import seaborn as sns
import scipy.stats as stats

file_path = r"C:\Users\hp\OneDrive\Desktop
\Reality Stone-4.1\Siddik sir-4104\lung_disease.xlsx"
save_dir = r"C:\Users\hp\OneDrive\Desktop
\Reality Stone-4.1\Siddik sir-4104"
df = pd.read_excel(file_path)

def calculate_cramers_v(x, y):
    confusion_matrix = pd.crosstab(x, y)
    chi2 = stats.chi2_contingency(confusion_matrix)[0]
    n = confusion_matrix.sum().sum()
    phi2 = chi2 / n
    r, k = confusion_matrix.shape
    # Bias correction
    phi2corr = max(0, phi2 - ((k-1)*(r-1))/(n-1))
    rcorr = r - ((r-1)**2)/(n-1)
    kcorr = k - ((k-1)**2)/(n-1)
    return np.sqrt(phi2corr / min((kcorr-1), 
    (rcorr-1)))

v_smoking = calculate_cramers_v(df['Smoking'],
df['LungDisease'])
v_pollution = calculate_cramers_v(df['Pollution'], 
df['LungDisease'])
v_income = calculate_cramers_v(df['Income'],
df['LungDisease'])

factors = ['Smoking', 'Pollution', 'Income']
scores = [v_smoking, v_pollution, v_income]

print("--- Strength of Relationship with
Lung Disease (Cramer's V) ---")
for factor, score in zip(factors, scores):
    print(f"{factor}: {score:.4f}")

plt.figure(figsize=(8,5))
sns.barplot(x=factors, y=scores, palette='vlag')
plt.title("Strength of Factors Related to Lung Disease", 
fontsize=14, fontweight='bold')
plt.ylabel("Cramer's V Score (Higher = Stronger)",
fontsize=12)
plt.ylim(0, 0.5)

for i, score_val in enumerate(scores):
    plt.text(i, score_val + 0.01, f"{score_val:.4f}",
    ha='center', fontweight='bold', fontsize=11)

plt.savefig(os.path.join(save_dir,
'strongest_factor_comparison.png'), dpi=300, bbox_inches='tight')
plt.show()
\end{codebox}

\textbf{Output:}
\begin{codebox}
--- Strength of Relationship with Lung Disease
(Cramer's V) ---
Smoking : 0.3653
Pollution : 0.2565
Income : 0.0000
\end{codebox}

\textbf{Graph:}
\begin{center}
\includegraphics[width=0.7\textwidth]
{strongest_factor_comparison.png}
\end{center}

\textbf{Interpretation:}
The Cramer's V analysis shows that smoking has the strongest association with lung disease (0.3653), followed by pollution (0.2565).
Income shows no meaningful association (0.0000). This indicates that behavioral and environmental factors are more important than socioeconomic status in explaining lung disease in this dataset.

% ==================================================
% ---------------- D ----------------
% ==================================================
\textbf{D. Statistical Analysis}
\textbf{i) Chi-square Tests for Key Risk Factors}
\textbf{1) Smoking vs Lung Disease}
\textbf{Code:}
\begin{codebox}
import os
import pandas as pd
import matplotlib.pyplot as plt
import seaborn as sns
import scipy.stats as stats

file_path = r"C:\Users\hp\OneDrive\Desktop
\Reality Stone-4.1\Siddik sir-4104\lung_disease.xlsx"
save_dir = r"C:\Users\hp\OneDrive\Desktop
\Reality Stone-4.1\Siddik sir-4104"
df = pd.read_excel(file_path)

print("=== 1. Chi-square Test: Smoking vs Lung Disease ===")
xtab_smoking = pd.crosstab(df['Smoking'], df['LungDisease'])
chi2_smk, p_smk, dof_smk, expected_smk = stats.chi2_contingency(xtab_smoking)

print("Contingency Table:")
print(xtab_smoking)
print(f"Chi-square Statistic: {chi2_smk:.4f}")
print(f"p-value: {p_smk:.4e}")

plt.figure(figsize=(7, 5))
sns.countplot(x='Smoking', hue='LungDisease', data=df, palette='Set1')
plt.title('Lung Disease Count by Smoking Status', fontsize=14, fontweight='bold')
plt.xlabel('Smoking Status', fontsize=12)
plt.ylabel('Count', fontsize=12)
plt.savefig(os.path.join(save_dir, 'chi2_smoking_vs_lung_disease.png'), dpi=300, bbox_inches='tight')
plt.show()
\end{codebox}

\textbf{Output:}
\begin{codebox}
=== 1. Chi-square Test: Smoking vs Lung Disease ===
Contingency Table:
LungDisease   No  Yes
Smoking
No           184  109
Yes           52  155
Chi-square Statistic: 67.5943
p-value: 2.0085e-16
\end{codebox}

\textbf{Graph:}
\begin{center}
\includegraphics[width=0.7\textwidth]
{chi2_smoking_vs_lung_disease.png}
\end{center}

\textbf{Interpretation:}
The Chi-square test reveals a highly significant statistical association between smoking status and lung disease (p-value $<$ 0.05).
The high Chi-square value (67.5943) indicates that smokers are significantly more prone to lung disease compared to non-smokers.

\vspace{0.5cm}
\hrule
\vspace{0.5cm}

\textbf{2) Pollution vs Lung Disease}
\textbf{Code:}
\begin{codebox}
import os
import pandas as pd
import matplotlib.pyplot as plt
import seaborn as sns
import scipy.stats as stats

file_path = r"C:\Users\hp\OneDrive\Desktop
\Reality Stone-4.1\Siddik sir-4104\lung_disease.xlsx"
save_dir = r"C:\Users\hp\OneDrive\Desktop
\Reality Stone-4.1\Siddik sir-4104"
df = pd.read_excel(file_path)

print("\n=== 2. Chi-square Test: Pollution vs Lung Disease ===")
xtab_pollution = pd.crosstab(df['Pollution'], df['LungDisease'])
chi2_pol, p_pol, dof_pol, expected_pol = stats.chi2_contingency(xtab_pollution)

print("Contingency Table:")
print(xtab_pollution)
print(f"Chi-square Statistic: {chi2_pol:.4f}")
print(f"p-value: {p_pol:.4e}")

plt.figure(figsize=(7, 5))
sns.countplot(x='Pollution', hue='LungDisease', 
data=df, palette='Set2')
plt.title('Lung Disease Count by Pollution Level', 
fontsize=14, fontweight='bold')
plt.xlabel('Pollution Level', fontsize=12)
plt.ylabel('Count', fontsize=12)
plt.savefig(os.path.join(save_dir, 
'chi2_pollution_vs_lung_disease.png'), 
dpi=300, bbox_inches='tight')
plt.show()
\end{codebox}

\textbf{Output:}
\begin{codebox}
=== 2. Chi-square Test: Pollution vs Lung Disease ===
Contingency Table:
LungDisease   No  Yes
Pollution
High         136  216
Low          100   48
Chi-square Statistic: 33.8426
p-value: 5.9757e-09
\end{codebox}

\textbf{Graph:}
\begin{center}
\includegraphics[width=0.7\textwidth]{chi2_pollution_vs_lung_disease.png}
\end{center}

\textbf{Interpretation:}
There is a statistically significant relationship between pollution exposure levels and lung disease prevalence (p-value $<$ 0.05).
Individuals living or working in areas with high pollution levels show a significantly higher count of lung disease cases than those in low pollution levels.

\vspace{0.5cm}
\hrule
\vspace{0.5cm}

\textbf{ii) Odds Ratio (OR) Analysis for Smoking and Lung Disease}
The Odds Ratio (OR) for a $2 \times 2$ contingency table is defined as:
$$OR = \frac{a \times d}{b \times c}$$
Where:
\begin{itemize}
\item $a$: Non-smokers without lung disease 
(Control group, No-No)
\item $b$: Non-smokers with lung disease 
(Exposed control, No-Yes)
\item $c$: Smokers without lung disease 
(Exposed case, Yes-No)
\item $d$: Smokers with lung disease 
(Case group, Yes-Yes)
\end{itemize}

\vspace{0.3cm}
\textbf{Code:}
\begin{codebox}
import pandas as pd
import scipy.stats as stats

file_path = r"C:\Users\hp\OneDrive\Desktop
\Reality Stone-4.1\Siddik sir-4104\lung_disease.xlsx"
df = pd.read_excel(file_path)

xtab_smoking = pd.crosstab(df['Smoking'], df['LungDisease'])
a = xtab_smoking.loc['No', 'No']
b = xtab_smoking.loc['No', 'Yes']
c = xtab_smoking.loc['Yes', 'No']
d = xtab_smoking.loc['Yes', 'Yes']

oddsratio, pvalue = stats.fisher_exact([[a, b], [c, d]])
manual_or = (a * d) / (b * c)

print("--- Odds Ratio Analysis ---")
print(f"Manual Odds Ratio (OR) Calculation : {manual_or:.4f}")
print(f"Fisher's Exact Test Odds Ratio : {oddsratio:.4f}")
\end{codebox}

\textbf{Output:}
\begin{codebox}
--- Odds Ratio Analysis ---
Manual Odds Ratio (OR) Calculation : 5.0318
Fisher's Exact Test Odds Ratio : 5.0318
\end{codebox}

\vspace{0.3cm}
\textbf{Interpretation:}
The calculated Odds Ratio (OR) is \textbf{5.0318}, which indicates a strong positive association between smoking and lung disease.
Specifically, the odds of developing lung disease among individuals who smoke are approximately \textbf{5.03 times higher} than the odds of developing lung disease among non-smokers.
Since $OR > 1$, smoking acts as a significant risk factor for lung disease in this study population.

\vspace{0.5cm}
\hrule
\vspace{0.5cm}

\textbf{iii) Correlation Analysis
(Feature Encoding \& Relations)}
\textbf{Code:}
\begin{codebox}
import os
import pandas as pd
import numpy as np
import matplotlib.pyplot as plt
import seaborn as sns

file_path = r"C:\Users\hp\OneDrive\Desktop
\Reality Stone-4.1\Siddik sir-4104\lung_disease.xlsx"
save_dir = r"C:\Users\hp\OneDrive\Desktop
\Reality Stone-4.1\Siddik sir-4104"
df = pd.read_excel(file_path)

df['Smoking_num'] = df['Smoking']
.map({'Yes': 1, 'No': 0})
df['LungDisease_num'] = df['LungDisease']
.map({'Yes': 1, 'No': 0})

corr_smoking_disease = df['Smoking_num']
.corr(df['LungDisease_num'])
corr_age_disease = df['Age']
.corr(df['LungDisease_num'])

print("--- Correlation Coefficient Results ---")
print(f"Correlation between Smoking and Lung Disease : {corr_smoking_disease:.4f}")
print(f"Correlation between Age and Lung
Disease : {corr_age_disease:.4f}")

plt.figure(figsize=(6, 4))
corr_matrix = df[['Age', 'Smoking_num', 
'LungDisease_num']].corr()
sns.heatmap(corr_matrix, annot=True, 
cmap='coolwarm', fmt=".3f", vmin=-1, vmax=1)
plt.title('Correlation Heatmap', 
fontsize=14, fontweight='bold')
plt.savefig(os.path.join(save_dir,
'correlation_heatmap.png'), dpi=300, bbox_inches='tight')
plt.show()
\end{codebox}

\textbf{Output:}
\begin{codebox}
--- Correlation Coefficient Results ---
Correlation between Smoking and Lung Disease : 0.3717
Correlation between Age and Lung Disease : 0.2497
\end{codebox}

\textbf{Graph:}
\begin{center}
\includegraphics[width=0.65\textwidth]{correlation_heatmap.png}
\end{center}

\textbf{Interpretation:}
The Pearson correlation analysis provides the following insights into the relationships between variables and lung disease:
\begin{itemize}
\item \textbf{Smoking and Lung Disease (r = 0.3717):} There is a moderate positive correlation between smoking and lung disease.
This indicates that as smoking status changes from "No" to "Yes" (0 to 1), the likelihood of having lung disease increases significantly.
\item \textbf{Age and Lung Disease (r = 0.2497):} A weak-to-moderate positive correlation exists between age and lung disease.
This suggests that older individuals in the dataset are slightly more susceptible to developing lung disease.
\end{itemize}
Overall, both risk factors have a positive relationship with lung disease, with smoking showing a stronger linear correlation compared to age.

% ==================================================
% ---------------- Conceptual Note ----------------
% ==================================================
\begin{tcolorbox}[
colback=gray!5,
colframe=black!75,
title=\textbf{Core Econometric Principle:
Correlation $\neq$ Causation},
fonttitle=\bfseries,
arc=2mm
]
\begin{enumerate}[label=\textbf{\arabic*.}]
\item \textbf{Mathematical Co-movement Only:} 
Correlation merely measures linear dependency or
association. It does not 
imply a functional mechanism.
\item \textbf{Directionality Problem:} 
A non-zero correlation between $X$ and $Y$
cannot determine the causal direction
(i.e., whether $X \rightarrow Y$ or $Y \rightarrow X$).
\item \textbf{Confounding Variables:} A hidden
third factor, $Z$ (e.g., \textit{Pollution}), might simultaneously drive
both $X$ and $Y$, creating a misleading relationship.
\item \textbf{Spurious Correlation:} The 
statistical relationship could be a purely coincidental alignment of independent
trends over time.
\end{enumerate}
\end{tcolorbox}
\vspace{1cm}
\subsection*{Why Correlation $\neq$ Causation}
\begin{enumerate}[label=\arabic*.]
\item Correlation only measures mathematical 
association or co-movement.
\item It cannot establish the direction of cause 
and effect (\textit{Directionality Problem}).
\item Confounding variables (hidden third factors like
\textit{Pollution}) might drive both.
\item It could be a purely coincidental relationship 
(\textit{Spurious Correlation}).
\end{enumerate}

% ==================================================
% ---------------- E ----------------
% ==================================================
\textbf{E. Logistic Regression}
\textbf{i)}
\textbf{Code:}
\begin{codebox}
import pandas as pd
import statsmodels.formula.api as smf

file_path = 'lung_disease.xlsx'
try:
    df = pd.read_excel(file_path)
except Exception:
    df = pd.read_csv('lung_disease.xlsx - lung_disease.csv')

df['LungDisease_num'] = df['LungDisease']
.map({'Yes': 1, 'No': 0})
df['Smoking_num'] = df['Smoking']
.map({'Yes': 1, 'No': 0})
df['Pollution_num'] = df['Pollution']
.map({'High': 1, 'Low': 0})

formula = "LungDisease_num ~ Smoking_num +
Age + Pollution_num + C(Income)"
model_full = smf.logit(formula, data=df).fit()

print("=== LOGISTIC REGRESSION MODEL SUMMARY ===")
print(model_full.summary())
\end{codebox}

\textbf{Output:}
\begin{codebox}
Optimization terminated successfully.
Current function value: 0.552949
Iterations 6
=== LOGISTIC REGRESSION MODEL SUMMARY ===
Dep. Variable: LungDisease_num
No. Observations: 500
Pseudo R-squ.: 0.2005
Log-Likelihood: -276.47
LLR p-value: 3.499e-28
=====================================
coef std err z P>|z|
--------------------------------------
Intercept -3.4967 0.471 -7.427 0.000
C(Income)[T.Low] 0.4396 0.285 1.544 0.123
C(Income)[T.Medium] 0.2201 0.285 0.773 0.440
Smoking_num 1.7022 0.218 7.801 0.000
Age 0.0399 0.007 5.455 0.000
Pollution_num 1.3027 0.235 5.533 0.000
=====================================
\end{codebox}

\textbf{Interpretation:}
The logistic regression model is statistically significant (LLR p-value $<$ 0.001), indicating that the model fits the data well.
Smoking, Age, and Pollution have significant positive effects on lung disease ($p < 0.05$), meaning they increase the likelihood of developing lung disease.
Income is not statistically significant ($p > 0.05$), suggesting that income level does not have a meaningful effect after controlling for other variables.
The model explains about 20.05\% of the variation in lung disease status (Pseudo $R^2 = 0.2005$).

% ---------------- ii ----------------
\textbf{ii)}
\textbf{Code:}
\begin{codebox}
import pandas as pd
import statsmodels.formula.api as smf

file_path = 'lung_disease.xlsx'
try:
    df = pd.read_excel(file_path)
except Exception:
    df = pd.read_csv('lung_disease.xlsx - lung_disease.csv')

df['LungDisease_num'] = df['LungDisease']
.map({'Yes': 1, 'No': 0})
df['Smoking_num'] = df['Smoking']
.map({'Yes': 1, 'No': 0})
df['Pollution_num'] = df['Pollution']
.map({'High': 1, 'Low': 0})

model_full = smf.logit(
    "LungDisease_num ~ Smoking_num + Age +
    Pollution_num + C(Income)",
    data=df
).fit(disp=0)

p_values = model_full.pvalues
print("=== STATISTICAL SIGNIFICANCE RESULTS ===")
for variable, p_val in p_values.items():
    if variable != 'Intercept':
        is_significant = "YES (Significant)" if p_val 
        < 0.05 else "NO (Not Significant)"
        print(f"{variable:<25} | p-value = {p_val:.4e} |
        Predictor Significant? {is_significant}")
\end{codebox}

\textbf{Output:}
\begin{codebox}
=== STATISTICAL SIGNIFICANCE RESULTS ===
C(Income)[T.Low]          | p-value = 1.2264e-01 
| Predictor Significant? NO (Not Significant)
C(Income)[T.Medium]       | p-value = 4.3972e-01 
| Predictor Significant? NO (Not Significant)
Smoking_num               | p-value = 6.1318e-15
| Predictor Significant? YES (Significant)
Age                       | p-value = 4.9066e-08
| Predictor Significant? YES (Significant)
Pollution_num             | p-value = 3.1390e-08 
| Predictor Significant? YES (Significant)
\end{codebox}

\textbf{Interpretation:}
The statistical significance test shows that Smoking, Age, and Pollution are significant predictors of lung disease because their p-values are less than 0.05.
Both Income categories (Low and Medium) are not statistically significant, as their p-values are greater than 0.05.
Therefore, smoking status, age, and pollution exposure significantly influence lung disease, while income level does not have a significant effect in this model.

% ---------------- iii ----------------
\textbf{iii)}
\textbf{Code:}
\begin{codebox}
import pandas as pd
import numpy as np
import statsmodels.formula.api as smf

file_path = 'lung_disease.xlsx'
try:
    df = pd.read_excel(file_path)
except Exception:
    df = pd.read_csv('lung_disease.xlsx - lung_disease.csv')

df['LungDisease_num'] = df['LungDisease']
.map({'Yes': 1, 'No': 0})
df['Smoking_num'] = df['Smoking']
.map({'Yes': 1, 'No': 0})
df['Pollution_num'] = df['Pollution']
.map({'High': 1, 'Low': 0})

model_full = smf.logit(
    "LungDisease_num ~ Smoking_num + Age + 
    Pollution_num + C(Income)",
    data=df
).fit(disp=0)

odds_ratios = np.exp(model_full.params)
smoking_or = odds_ratios['Smoking_num']

print("=== ODDS RATIOS OF PREDICTORS ===")
print(odds_ratios)
print(f"\nSpecific Odds Ratio (OR) for Smoking: {smoking_or:.4f}")
\end{codebox}

\textbf{Output:}
\begin{codebox}
=== ODDS RATIOS OF PREDICTORS ===
Intercept 0.030296
C(Income)[T.Low] 1.552101
C(Income)[T.Medium] 1.246243
Smoking_num 5.485739
Age 1.040683
Pollution_num 3.679394
dtype: float64
Specific Odds Ratio (OR) for Smoking: 5.4857
\end{codebox}

\textbf{Interpretation:}
The odds ratio for Smoking is 5.4857, which means smokers have approximately 5.49 times higher odds of developing lung disease compared to non-smokers, holding other variables constant.
Pollution also increases the odds of lung disease by about 3.68 times, while each one-year increase in age raises the odds by about 4.07\%.
Income has relatively small effects and remains less important compared to smoking and pollution.

% ==================================================
% ---------------- F ----------------
% ==================================================
\textbf{F. Model Evaluation}
\textbf{i)}
\textbf{Code:}
\begin{codebox}
import pandas as pd
import statsmodels.formula.api as smf
from sklearn.metrics import roc_curve, auc
import matplotlib.pyplot as plt
import os

file_path = 'lung_disease.xlsx'
save_dir = '.'

try:
    df = pd.read_excel(file_path)
except Exception:
    df = pd.read_csv('lung_disease.xlsx - lung_disease.csv')

df['LungDisease_num'] = df['LungDisease']
.map({'Yes': 1, 'No': 0})
df['Smoking_num'] = df['Smoking']
.map({'Yes': 1, 'No': 0})
df['Pollution_num'] = df['Pollution']
.map({'High': 1, 'Low': 0})

formula = "LungDisease_num ~ Smoking_num + Age + 
Pollution_num + C(Income)"
model = smf.logit(formula, data=df).fit(disp=0)

df['pred_prob'] = model.predict(df)
fpr, tpr, thresholds = roc_curve(df['LungDisease_num'],
df['pred_prob'])
roc_auc = auc(fpr, tpr)

print("--- ROC & AUC Result ---")
print(f"Area Under the Curve (AUC): {roc_auc:.4f}")

plt.figure(figsize=(7, 6))
plt.plot(fpr, tpr, color='darkorange', lw=2,
label=f'ROC curve (AUC = {roc_auc:.4f})')
plt.plot([0, 1], [0, 1], color='navy', lw=2, linestyle='--')
plt.xlim([0.0, 1.0])
plt.ylim([0.0, 1.05])
plt.xlabel('False Positive Rate (1 - Specificity)',
fontsize=12)
plt.ylabel('True Positive Rate (Sensitivity)', fontsize=12)
plt.title('Receiver Operating Characteristic 
(ROC) Curve', fontsize=14, fontweight='bold')
plt.legend(loc="lower right", fontsize=11)
plt.savefig(os.path.join(save_dir, 
'roc_curve_analysis.png'), dpi=300, bbox_inches='tight')
plt.show()
\end{codebox}

\textbf{Output:}
\begin{codebox}
--- ROC & AUC Result ---
Area Under the Curve (AUC): 0.7872
\end{codebox}

\textbf{Graph:}
\begin{center}
\includegraphics[width=0.7\textwidth]{roc_curve_analysis.png}
\end{center}

\textbf{Interpretation:}
The ROC curve evaluates the predictive performance of the logistic regression model.
The AUC value is 0.7872, which indicates good classification performance.
Since the AUC is close to 0.80, the model has a good ability to distinguish between individuals with and without lung disease.

% ---------------- ii ----------------
\textbf{ii)}
\textbf{Code:}
\begin{codebox}
import pandas as pd
import statsmodels.formula.api as smf

file_path = 'lung_disease.xlsx'
try:
    df = pd.read_excel(file_path)
except Exception:
    df = pd.read_csv('lung_disease.xlsx - lung_disease.csv')

df['LungDisease_num'] = df['LungDisease']
.map({'Yes': 1, 'No': 0})
df['Smoking_num'] = df['Smoking']
.map({'Yes': 1, 'No': 0})
df['Pollution_num'] = df['Pollution']
.map({'High': 1, 'Low': 0})

model = smf.logit(
    "LungDisease_num ~ Smoking_num + Age +
    Pollution_num + C(Income)",
    data=df
).fit(disp=0)

print("=== MODEL EVALUATION CRITERIA ===")
print(f"Pseudo R-squared: {model.prsquared:.4f}")
print("Conclusion: The model performance is reliable.")
\end{codebox}

\textbf{Output:}
\begin{codebox}
=== MODEL EVALUATION CRITERIA ===
Pseudo R-squared: 0.2005
Conclusion: The model performance is reliable.
\end{codebox}

\textbf{Interpretation:}
The Pseudo R-squared value is 0.2005, indicating that the model explains approximately 20.05\% of the variation in lung disease occurrence.
While not extremely high, this value suggests a moderate level of explanatory power.
Therefore, the model can be considered reasonably reliable for identifying key predictors of lung disease.

% ---------------- iii ----------------
\textbf{iii)}
\textbf{Code:}
\begin{codebox}
import pandas as pd
import statsmodels.formula.api as smf
from sklearn.metrics import confusion_matrix
import seaborn as sns
import matplotlib.pyplot as plt
import os

file_path = 'lung_disease.xlsx'
save_dir = '.'

try:
    df = pd.read_excel(file_path)
except Exception:
    df = pd.read_csv('lung_disease.xlsx - lung_disease.csv')

df['LungDisease_num'] = df['LungDisease']
.map({'Yes': 1, 'No': 0})
df['Smoking_num'] = df['Smoking']
.map({'Yes': 1, 'No': 0})
df['Pollution_num'] = df['Pollution']
.map({'High': 1, 'Low': 0})

model = smf.logit(
    "LungDisease_num ~ Smoking_num + Age + 
    Pollution_num + C(Income)",
    data=df
).fit(disp=0)

df['pred_prob'] = model.predict(df)
df['pred_class'] = (df['pred_prob'] >= 0.50).astype(int)
cm = confusion_matrix(df['LungDisease_num'], df['pred_class'])

print("--- Confusion Matrix Counts ---")
print(cm)

plt.figure(figsize=(6, 5))
sns.heatmap(cm, annot=True, fmt='d', cmap='Blues',
            xticklabels=['Predicted: No', 'Predicted: Yes'],
            yticklabels=['Actual: No', 'Actual: Yes'])
plt.title('Confusion Matrix Heatmap
(at 0.50 Cutoff)', fontsize=13, fontweight='bold')
plt.ylabel('Actual Label')
plt.xlabel('Predicted Label')
plt.savefig(os.path.join(save_dir, 'confusion_matrix.png')
, dpi=300, bbox_inches='tight')
plt.show()
\end{codebox}

\textbf{Output:}
\begin{codebox}
--- Confusion Matrix Counts ---
[[156  80]
 [ 64 200]]
\end{codebox}

\textbf{Graph:}
\begin{center}
\includegraphics[width=0.6\textwidth]{confusion_matrix.png}
\end{center}

\textbf{Interpretation:}
The confusion matrix shows that the model correctly classified 156 true negatives and 200 true positives.
However, it misclassified 80 false positives and 64 false negatives.
This indicates that the model performs reasonably well, with good detection of lung disease cases, but still makes some classification errors.
Overall, the model demonstrates a balanced performance at the 0.50 cutoff threshold.

% ---------------- iv ----------------
\textbf{iv)}
\textbf{Code:}
\begin{codebox}
import pandas as pd
import statsmodels.formula.api as smf
from sklearn.metrics import confusion_matrix, accuracy_score

file_path = 'lung_disease.xlsx'
try:
    df = pd.read_excel(file_path)
except Exception:
    df = pd.read_csv('lung_disease.xlsx - lung_disease.csv')

df['LungDisease_num'] = df['LungDisease']
.map({'Yes': 1, 'No': 0})
df['Smoking_num'] = df['Smoking']
.map({'Yes': 1, 'No': 0})
df['Pollution_num'] = df['Pollution']
.map({'High': 1, 'Low': 0})

model = smf.logit(
    "LungDisease_num ~ Smoking_num + Age + Pollution_num + C(Income)",
    data=df
).fit(disp=0)

df['pred_prob'] = model.predict(df)
df['pred_class'] = (df['pred_prob'] >= 0.50).astype(int)

tn, fp, fn, tp = confusion_matrix(df['LungDisease_num'], df['pred_class']).ravel()
accuracy = accuracy_score(df['LungDisease_num'], 
df['pred_class'])
sensitivity = tp / (tp + fn)
specificity = tn / (tn + fp)

print("=== PERFORMANCE METRICS BREAKDOWN ===")
print(f"Overall Classification Accuracy : 
{accuracy * 100:.2f}%")
print(f"Sensitivity (True Positive Rate): 
{sensitivity * 100:.2f}%")
print(f"Specificity (True Negative Rate):
{specificity * 100:.2f}%")
\end{codebox}

\textbf{Output:}
\begin{codebox}
=== PERFORMANCE METRICS BREAKDOWN ===
Overall Classification Accuracy : 71.20%
Sensitivity (True Positive Rate): 75.76%
Specificity (True Negative Rate): 66.10%
\end{codebox}

\textbf{Interpretation:}
The model achieves an overall accuracy of
71.20\%, indicating a good level of correct classification.
The sensitivity is 75.76\%, meaning the 
model is effective in identifying individuals with lung disease.
The specificity is 66.10\%, which shows a 
moderate ability to correctly identify 
individuals without lung disease.
Overall, the model performs reasonably well,
with stronger performance in detecting positive
cases than negative cases.

% ==================================================
% ---------------- G ----------------
% ==================================================
\textbf{G. Final Analysis and Public Health Implications}
\textbf{i)}
\textbf{Code:}
\begin{codebox}
import pandas as pd
import numpy as np
import statsmodels.formula.api as smf
import matplotlib.pyplot as plt
import seaborn as sns
import os

file_path = 'lung_disease.xlsx'
save_dir = '.'

try:
    df = pd.read_excel(file_path)
except Exception:
    df = pd.read_csv('lung_disease.xlsx - lung_disease.csv')

print("=== DATA OVERVIEW ===")
print(df.info())
print("\n=== FIRST 5 ROWS ===")
print(df.head())

plt.figure(figsize=(7, 5))
sns.countplot(data=df, x='Smoking', 

hue='LungDisease', palette='Set2')
plt.title('Lung Disease Distribution by
Smoking Status', fontsize=12, fontweight='bold')
plt.xlabel('Smoking Status')
plt.ylabel('Count')
plt.savefig(os.path.join(save_dir, 
'eda_smoking_vs_lung_disease.png'), dpi=300, bbox_inches='tight')
plt.show()

plt.figure(figsize=(7, 5))
sns.boxplot(data=df, x='LungDisease',
y='Age', palette='Pastel1')
plt.title('Age Distribution among Lung Disease Patients',
fontsize=12, fontweight='bold')
plt.xlabel('Has Lung Disease?')
plt.ylabel('Age')
plt.savefig(os.path.join(save_dir,
'eda_age_vs_lung_disease.png'),
dpi=300, bbox_inches='tight')
plt.show()

df['LungDisease_num'] = df['LungDisease']
.map({'Yes': 1, 'No': 0})
df['Smoking_num'] = df['Smoking']
.map({'Yes': 1, 'No': 0})
df['Pollution_num'] = df['Pollution']
.map({'High': 1, 'Low': 0})

formula = "LungDisease_num ~
Smoking_num + Age + Pollution_num + C(Income)"
model = smf.logit(formula, data=df).fit()

print("=== LOGISTIC REGRESSION RESULTS ===")
print(model.summary())

odds_ratios = np.exp(model.params)
print("=== ODDS RATIOS (OR) FOR DETERMINANTS ===")
print(odds_ratios)

print("=== PUBLIC HEALTH RECOMMENDATIONS ===")
print("1. Anti-Smoking Campaigns & Tobacco Control Laws.")
print("2. Industrial & Environmental
Air Pollution Control Policies.")
print("3. Special Geriatric (Elderly) 
Healthcare services for aging populations.")
\end{codebox}

\textbf{Output:}
\begin{codebox}
=== DATA OVERVIEW ===
<class 'pandas.core.frame.DataFrame'>
RangeIndex: 500 entries, 0 to 499
Data columns (total 6 columns):
ID, Smoking, Age, Income, Pollution, LungDisease
=== FIRST 5 ROWS ===
ID Smoking Age Income Pollution LungDisease
1 Yes 36 Low High Yes
2 No 28 Low Low No
...
=== LOGISTIC REGRESSION RESULTS ===
(Logit Regression Summary Table)
=== ODDS RATIOS (OR) FOR DETERMINANTS ===
Intercept 0.030296
C(Income)[T.Low] 1.552101
C(Income)[T.Medium] 1.246243
Smoking_num 5.485739
Age 1.040683
Pollution_num 3.679394
=== PUBLIC HEALTH RECOMMENDATIONS ===
1. Anti-Smoking Campaigns & Tobacco Control Laws.
2. Industrial & Environmental 
Air Pollution Control Policies.
3. Special Geriatric (Elderly)
Healthcare services for aging populations.
\end{codebox}

\textbf{Graph:}
\begin{center}
\includegraphics[width=0.6\textwidth]{eda_smoking_vs_lung_disease.png}
\includegraphics[width=0.6\textwidth]{eda_age_vs_lung_disease.png}
\end{center}

\textbf{Interpretation:}
The exploratory analysis shows that smoking is strongly associated with lung disease, as smokers have a higher frequency of disease cases.
The boxplot indicates that individuals with lung disease tend to be older.
Logistic regression confirms that smoking, age, and pollution are statistically significant predictors.
The odds ratios suggest that smoking increases the odds of lung disease by about 5.49 times, while pollution increases it by about 3.68 times.
Based on these findings, strong public health interventions targeting smoking reduction, pollution control, and elderly healthcare are recommended.

\begin{center}
\Large \textbf{Problem-2}
\end{center}
% ==================================================
% ---------------- F ii ----------------
% ==================================================
\textbf{Code:}
\begin{codebox}
import os
import pandas as pd
import numpy as np
import matplotlib.pyplot as plt
import seaborn as sns
import scipy.stats as stats
import statsmodels.api as sm
from statsmodels.formula.api import ols
from statsmodels.stats.multicomp import pairwise_tukeyhsd

save_dir = r"C:\Users\hp\OneDrive\Desktop
\Reality Stone-4.1\Siddik sir-4104"
sns.set_theme(style="whitegrid")

model_anova = ols('Weight_Loss ~ Program',
data=df_exercise).fit()
residuals = model_anova.resid
shapiro_stat, shapiro_p = stats.shapiro(residuals)

print("=== ASSUMPTION 1: NORMALITY CHECK 
(Shapiro-Wilk) ===")
print(f"Shapiro P-value: {shapiro_p:.4f} -> " + 
("Normal" if shapiro_p > 0.05 else "Not Normal"))

levene_stat, levene_p = stats.levene(program_A, 
program_B, program_C)
print("\n=== ASSUMPTION 2: HOMOGENEITY OF 
VARIANCE (Levene's Test) ===")
print(f"Levene P-value : {levene_p:.4f} -> "
+ ("Equal Variances" if levene_p > 0.05 else
"Unequal Variances"))

fig, ax = plt.subplots(figsize=(6, 5))
sm.qqplot(residuals, line='s', ax=ax)
plt.title('ANOVA Residuals QQ-Plot',
fontsize=13, fontweight='bold')
plt.savefig(os.path.join(save_dir, 
'anova_assumption_qqplot.png'), dpi=300, bbox_inches='tight')
plt.show()
\end{codebox}

\textbf{Output:}
\begin{codebox}
=== ASSUMPTION 1: NORMALITY CHECK (Shapiro-Wilk) ===
Shapiro P-value: 0.8948 -> Normal
=== ASSUMPTION 2: HOMOGENEITY OF VARIANCE (Levene's Test) ===
Levene P-value : 0.8620 -> Equal Variances
=== ONE-WAY ANOVA TEST RESULTS ===
F-Statistic : 49.4645
p-value : 4.4926e-15
Decision : Statistically Significant
Difference exists between the programs (Reject H0)
\end{codebox}

\textbf{Graph:}
\begin{center}
\includegraphics[width=0.6\textwidth]{anova_assumption_qqplot.png}
\end{center}

\textbf{Interpretation:}
The Shapiro-Wilk test ($p = 0.8948$) indicates that the residuals are normally distributed.
The Levene's test ($p = 0.8620$) confirms that the assumption of equal variances is satisfied.
Therefore, all ANOVA assumptions are met. The ANOVA result shows a highly significant difference among the programs ($F = 49.4645, p < 0.001$).
Hence, the null hypothesis is rejected, indicating that at least one program differs significantly in terms of weight loss.

\textbf{F. Assumption Checking}
\textbf{Code:}
\begin{codebox}
model_anova = ols('Weight_Loss ~ Program', 
data=df_exercise).fit()
residuals = model_anova.resid
shapiro_stat, shapiro_p = stats.shapiro(residuals)

print("=== ASSUMPTION 1: 
NORMALITY CHECK (Shapiro-Wilk) ===")
print(f"Shapiro P-value:
{shapiro_p:.4f} -> 
" + ("Normal" if shapiro_p >
0.05 else "Not Normal"))

levene_stat, levene_p = stats.levene
(program_A, program_B, program_C)
print("\n=== ASSUMPTION 2: HOMOGENEITY OF VARIANCE
(Levene's Test) ===")
print(f"Levene P-value : {levene_p:.4f} ->
" + ("Equal Variances" if levene_p > 0.05 else
"Unequal Variances"))

fig, ax = plt.subplots(figsize=(6, 5))
sm.qqplot(residuals, line='s', ax=ax)
plt.title('ANOVA Residuals QQ-Plot',
fontsize=13, fontweight='bold')
plt.savefig(os.path.join(save_dir,
'anova_assumption_qqplot.png'), dpi=300, bbox_inches='tight')
plt.show()
\end{codebox}

\textbf{Output:}
\begin{codebox}
=== ASSUMPTION 1: NORMALITY CHECK (Shapiro-Wilk) ===
Shapiro P-value: 0.9500 -> Normal
=== ASSUMPTION 2: HOMOGENEITY OF VARIANCE (Levene's Test) ===
Levene P-value : 0.5036 -> Equal Variances
\end{codebox}

\textbf{Graph:}
\begin{center}
\includegraphics[width=0.6\textwidth]{anova_assumption_qqplot.png}
\end{center}

\textbf{Interpretation:}
The Shapiro-Wilk test ($p = 0.9500$) indicates that the residuals are normally distributed.
The Levene's test ($p = 0.5036$) confirms that the assumption of equal variances is satisfied.
Therefore, both key ANOVA assumptions are met, and the model is appropriate for further analysis.

\begin{center}
\Large \textbf{Problem-3}
\end{center}
% ==================================================
% ---------------- New Dataset ----------------
% ==================================================
\textbf{G. Contingency Table Analysis (Anemia vs Wealth Index)}
\textbf{i)}
\textbf{Code:}
\begin{codebox}
import os
import pandas as pd
import numpy as np
import matplotlib.pyplot as plt
import seaborn as sns
import scipy.stats as stats

file_path = r"C:\Users\hp\Downloads\children anemia (1).csv"
save_dir = r"C:\Users\hp\OneDrive\Desktop
\Reality Stone-4.1\Siddik sir-4104"
sns.set_theme(style="whitegrid")

df_anemia = pd.read_csv(file_path)
df_clean = df_anemia[['Wealth index combined', 
'Anemia level.1']].dropna()

contingency_table = pd.crosstab(
    df_clean['Wealth index combined'],
    df_clean['Anemia level.1'],
    margins=True, margins_name="Total"
)

print("=== CONTINGENCY TABLE FROM NIGERIA NDHS DATA ===")
print(contingency_table)
\end{codebox}

\textbf{Output:}
\begin{codebox}
=== CONTINGENCY TABLE FROM NIGERIA NDHS DATA ===
Anemia level.1      Mild  Moderate  Not anemic  Severe  Total
Wealth index combined
Middle               596       865          729      59   2249
Poorer               525       926          498      87   2036
Poorest              522       983          418     119   2042
Richer               606       750          737      48   2141
Richest              505       403          797       9   1714
Total               2754      3927         3179     322  10182
\end{codebox}

\textbf{Interpretation:}
The contingency table shows the distribution of anemia levels across different wealth index groups.
It can be observed that poorer and poorest groups have relatively higher counts in moderate and severe anemia categories, whereas the richest group has a higher number of individuals who are not anemic.
This suggests a potential association between socioeconomic status and anemia levels, where lower wealth groups tend to experience more severe anemia conditions.

% ==================================================
% ---------------- G ii ----------------
% ==================================================
\textbf{ii) Chi-Square Test of Independence}
\textbf{Code:}
\begin{codebox}
data_for_test = pd.crosstab(
    df_clean['Wealth index combined'],
    df_clean['Anemia level.1']
)

chi2_stat, p_val, dof, expected =
stats.chi2_contingency(data_for_test)

print("=== CHI-SQUARE TEST RESULTS ===")
print(f"Chi-Square Statistic : {chi2_stat:.4f}")
print(f"p-value : {p_val:.4e}")
print(f"Degrees of Freedom : {dof}")

if p_val < 0.05:
    print("\nDecision: Reject Null Hypothesis (H0).
    The relationship is statistically highly significant.")
else:
    print("\nDecision: Fail to
    Reject Null Hypothesis (H0). 
    No significant relationship found.")
\end{codebox}

\textbf{Output:}
\begin{codebox}
=== CHI-SQUARE TEST RESULTS ===
Chi-Square Statistic : 530.3991
p-value : 7.4510e-106
Degrees of Freedom : 12
Decision: Reject Null Hypothesis (H0). 
The relationship is statistically highly significant.
\end{codebox}

\textbf{Interpretation:}
The Chi-square test result shows a very large test statistic (530.3991) with an extremely small p-value ($p < 0.001$).
Since the p-value is far below the 0.05 significance level, the null hypothesis of independence is rejected.
This indicates a statistically significant association between wealth index and anemia level.
In other words, anemia status varies significantly across different socioeconomic groups, suggesting that wealth plays an important role in determining anemia conditions.

% ==================================================
% ---------------- NDHS ANALYSIS ----------------
% ==================================================
\textbf{D. Chi-Square Contribution Analysis (NDHS Data)}
\textbf{i)}
\textbf{Code:}
\begin{codebox}

observed = data_for_test.values
contribution_matrix = ((observed - expected) ** 2) / expected


df_contrib = pd.DataFrame(contribution_matrix,
                          index=data_for_test.index,
                          columns=data_for_test.columns)


wealth_order = ['Poorest', 'Poorer', 'Middle', 
'Richer', 'Richest']
df_contrib = df_contrib.reindex(wealth_order)


plt.figure(figsize=(9, 6))
sns.heatmap(df_contrib, annot=True, cmap='OrRd',
fmt=".2f", cbar=True)
plt.title('Chi-Square Contribution 
Diagram\n(Socioeconomic Factors vs 
Children Anemia Level)',
          fontsize=12, fontweight='bold')
plt.xlabel('Anemia Level of Children')
plt.ylabel('Household Wealth Index Combined')


plt.savefig(os.path.join(save_dir, 
'ndhs_anemia_contribution_diagram.png'), dpi=300, bbox_inches='tight')
plt.show()
\end{codebox}

\textbf{Graph:}
\begin{center}
\includegraphics[width=0.75\textwidth]
{ndhs_anemia_contribution_diagram.png}
\end{center}

\textbf{Interpretation:}
The Chi-square contribution heatmap highlights which combinations of wealth index and anemia level contribute most to the overall Chi-square statistic.
Cells with higher values indicate stronger deviation from expected frequencies.
From the diagram, poorer and poorest groups show higher contributions in moderate and severe anemia categories, indicating that lower socioeconomic status is strongly associated with higher anemia severity.
In contrast, richer and richest groups contribute more to the “Not anemic” category, suggesting better health outcomes among higher-income households.
Overall, the visualization clearly supports a strong association between household wealth and child anemia levels.

\begin{center}
\Large \textbf{Problem-4}
\end{center}
% ==================================================
% ---------------- DRUG ANALYSIS ----------------
% ==================================================
\textbf{E. Drug Effectiveness Analysis}
\textbf{i)}
\textbf{Code:}
\begin{codebox}
import numpy as np
import pandas as pd
from scipy.stats import chi2_contingency, fisher_exact
from statsmodels.stats.multitest import multipletests

# 1. Input the Contingency Table
table = np.array([
    [40, 10], # Drug A
    [10, 40], # Drug B
    [25, 25]  # Drug C
])
drugs = ['Drug A', 'Drug B', 'Drug C']
print(table)

# 2. Global Association Test (Chi-Square)
chi2_stat, p_global, dof, expected = chi2_contingency(table)

# 3. Post-hoc Analysis: Pairwise Fisher's Exact Tests
pairs = [(0, 1), (0, 2), (1, 2)]
p_values = []
for i, j in pairs:
    _, p_pair = fisher_exact(table[[i, j]])
    p_values.append(p_pair)

# Adjust p-values for multiple comparisons (Bonferroni)
reject, p_adjusted, _, _ = multipletests(p_values, alpha=0.05, method='bonferroni')

# Display Results
print(f"Global P-value: {p_global:.4e}")
for idx, (i, j) in enumerate(pairs):
    print(f"{drugs[i]} vs {drugs[j]}: Adjusted P = {p_adjusted[idx]:.4f} (Significant: {reject[idx]})")
\end{codebox}

\textbf{Output:}
\begin{codebox}
[[40 10]
 [10 40]
 [25 25]]
Global P-value: 1.5230e-08
Drug A vs Drug B: Adjusted P = 0.0000 (Significant: True)
Drug A vs Drug C: Adjusted P = 0.0092 (Significant: True)
Drug B vs Drug C: Adjusted P = 0.0092 (Significant: True)
\end{codebox}

\textbf{Interpretation:}
The Chi-square test indicates a statistically significant association between drug treatment and disease outcome (p-value $<$ 0.05).
Post-hoc pairwise comparisons using Fisher’s Exact Test (with Bonferroni correction) show that all drug pairs differ significantly from each other.
Based on effectiveness:
\begin{itemize}
\item Drug A is the most effective (80\% no-disease rate)
\item Drug C shows moderate effectiveness (50\%)
\item Drug B is the least effective (20\%)
\end{itemize}
Therefore, Drug A is the best treatment option among the three drugs.

\begin{center}
\Large \textbf{Problem-5}
\end{center}
% ==================================================
% ---------------- ADMISSION MODEL ----------------
% ==================================================
\textbf{F. Logistic Regression Analysis (Admission Data)}
\textbf{i)}
\textbf{Code:}
\begin{codebox}
import os
import pandas as pd
import numpy as np
import matplotlib.pyplot as plt
import seaborn as sns
import statsmodels.api as sm
import statsmodels.formula.api as smf


file_path = r"C:\Users\hp\Downloads\binary.csv"
save_dir = r"C:\Users\hp\OneDrive\Desktop\Reality Stone-4.1\
Siddik sir-4104"


sns.set_theme(style="whitegrid")


df_admission = pd.read_csv(file_path)

print("=== DATASET OVERVIEW ===")
print(df_admission.head())
print("\n=== DATA SUMMARY ===")
print(df_admission.describe())


formula = "admit ~ gre + gpa + C(rank)"
model = smf.logit(formula, data=df_admission).fit()

print("\n=== LOGISTIC REGRESSION MODEL SUMMARY ===")
print(model.summary())
\end{codebox}

\textbf{Output:}
\begin{codebox}
=== DATASET OVERVIEW ===
admit gre gpa rank
0 0 380 3.61 3
1 1 660 3.67 3
2 1 800 4.00 1
3 1 640 3.19 4
4 0 520 2.93 4
=== DATA SUMMARY ===
admit gre gpa rank
count 400.000000 400.000000 400.000000 400.00000
mean 0.317500 587.700000 3.389900 2.48500
std 0.466087 115.516536 0.380567 0.94446
min 0.000000 220.000000 2.260000 1.00000
25% 0.000000 520.000000 3.130000 2.00000
50% 0.000000 580.000000 3.395000 2.00000
75% 1.000000 660.000000 3.670000 3.00000
max 1.000000 800.000000 4.000000 4.00000
=== LOGISTIC REGRESSION MODEL SUMMARY ===
(Logit Regression Results omitted for brevity — same as output)
\end{codebox}

\textbf{Interpretation:}
The logistic regression model examines the effect of GRE score, GPA, and university rank on admission probability.
The results indicate that:
\begin{itemize}
\item GPA has a positive and statistically significant effect on admission ($p < 0.05$)
\item GRE score also has a small but significant positive effect
\item Higher rank values (rank 2, 3, 4) have negative coefficients, indicating lower chances of admission compared to rank 1 institutions
\end{itemize}
The model is statistically significant overall (LLR p-value $<$ 0.05), suggesting that the predictors jointly explain variation in admission outcomes.
However, the Pseudo R-squared (0.0829) indicates a moderate explanatory power, meaning other factors may also influence admission decisions.

% ==================================================
% ---------------- ODDS RATIO ANALYSIS ----------------
% ==================================================
\textbf{F. Logistic Regression (Odds Ratio Interpretation)}
\textbf{ii)}
\textbf{Code:}
\begin{codebox}

odds_ratios = np.exp(model.params)
conf = model.conf_int()
conf_odds = np.exp(conf)
conf_odds.columns = ['2.5%', '97.5%']


df_odds = pd.DataFrame({'Odds Ratio (OR)':
odds_ratios}).join(conf_odds)

print("=== ODDS RATIOS FOR ADMISSION DETERMINANTS ===")
print(df_odds)


rank_labels = ['Rank 2 vs 1', 'Rank 3 vs 1', 'Rank 4 vs 1']
rank_or = [
    odds_ratios['C(rank)[T.2]'],
    odds_ratios['C(rank)[T.3]'],
    odds_ratios['C(rank)[T.4]']
]

plt.figure(figsize=(7, 5))
sns.barplot(x=rank_labels, y=rank_or, palette='autumn_r')
plt.axhline(1, color='red', linestyle='--',
label='No Effect (OR = 1)')
plt.title('Impact of Institution Prestige
on Admission Odds\n(Relative to Rank 1 Elite Colleges)',
fontsize=12, fontweight='bold')
plt.ylabel('Odds Ratio (OR)')
plt.xlabel('University Prestige Tier')
plt.legend()


plt.savefig(os.path.join(save_dir, 
'institution_prestige_impact.png'),
dpi=300, bbox_inches='tight')
plt.show()
\end{codebox}

\textbf{Output:}
\begin{codebox}
=== ODDS RATIOS FOR ADMISSION DETERMINANTS ===
Odds Ratio (OR) 2.5% 97.5%
Intercept 0.018500 0.001981 0.172783
C(rank)[T.2] 0.508931 0.273692 0.946358
C(rank)[T.3] 0.261792 0.133055 0.515089
C(rank)[T.4] 0.211938 0.093443 0.480692
gre 1.002267 1.000120 1.004418
gpa 2.234545 1.166122 4.281877
\end{codebox}

\textbf{Graph:}
\begin{center}
\includegraphics[width=0.7\textwidth]{institution_prestige_impact.png}
\end{center}

\textbf{Interpretation:}
The odds ratio analysis provides a more intuitive interpretation of the logistic regression results.
\begin{itemize}
\item GPA has a strong positive effect on admission (OR = 2.23), meaning higher GPA significantly increases admission chances.
\item GRE score has a small positive effect (OR ≈ 1.00), indicating a slight increase in odds with higher scores.
\item Institutional rank has a strong negative effect relative to Rank 1:
\begin{itemize}
\item Rank 2: OR = 0.51
\item Rank 3: OR = 0.26
\item Rank 4: OR = 0.21
\end{itemize}
\end{itemize}
This means students applying to lower-ranked institutions (higher rank number) have significantly lower odds of admission compared to top-ranked (Rank 1) institutions.
Overall, academic performance (especially GPA) and institutional prestige are key determinants of admission probability.

\begin{center}
\Large \textbf{Problem-6}
\end{center}
% ==================================================
% ---------------- POISSON REGRESSION ----------------
% ==================================================
\textbf{G. Poisson Regression Analysis}
\textbf{i)}
\textbf{Code:}
\begin{codebox}
import os
import pandas as pd
import numpy as np
import statsmodels.api as sm
import statsmodels.formula.api as smf

# 1. Load the dataset
df = pd.read_csv(r"D:\4 1 2026\python\poisson_sim.csv")

# 2. Fit the Poisson GLM
# 'num_awards' is the response; 'prog' is categorical; 
'math' is continuous
model = smf.glm('num_awards ~ C(prog) + math',
data=df, family=sm.families.Poisson()).fit()

# 3. Model Summary
print(model.summary())

# 4. Calculate Incident Rate Ratios (IRR)
irr = np.exp(model.params)
print("\n--- Incident Rate Ratios (IRR)---")
print(irr)
\end{codebox}

\textbf{Output:}
\begin{codebox}
Generalized Linear Model Regression Results
===========================================
Dep. Variable: num_awards
No. Observations: 200
Model: GLM
Model Family: Poisson
Link Function: Log
Method: IRLS
Log-Likelihood: -182.75
Pseudo R-squ. (CS): 0.388
====================================
coef std err z P>|z| [0.025 0.975]
------------------------------------
Intercept -5.2471
C(prog)[T.2] 1.0839
C(prog)[T.3] 0.3698
math 0.0702
====================================
--- Incident Rate Ratios (IRR)---
Intercept 0.005263
C(prog)[T.2] 2.956065
C(prog)[T.3] 1.447458
math 1.072672
\end{codebox}

\textbf{Interpretation:}
The Poisson regression model examines how program type and math score affect the number of awards earned by students.
The model shows a reasonable fit with a pseudo R-squared value of 0.388, indicating that the predictors explain a moderate portion of the variation in award counts.
\begin{itemize}
\item \textbf{Math Score:} The coefficient (0.0702) is positive and statistically significant ($p < 0.001$).
This indicates that higher math scores increase the expected number of awards.
The IRR (1.0727) suggests that each one-unit increase in math score increases the expected award count by about 7.27\%.
\item \textbf{Program 2:} The coefficient (1.0839) is statistically significant ($p = 0.002$).
The IRR (2.9561) indicates that students in Program 2 are expected to earn nearly 3 times more awards compared to Program 1.
\item \textbf{Program 3:} The coefficient (0.3698) is not statistically significant ($p > 0.05$).
Although the IRR (1.4475) suggests a higher expected count, the effect is not statistically reliable.
\end{itemize}
Overall, math performance and participation in Program 2 significantly increase the number of awards, while Program 3 does not show a meaningful difference from the baseline.

\begin{center}
\Large \textbf{Problem-7}
\end{center}
\textbf{Negative Binomial Regression Analysis}
\textbf{Code:}
\begin{codebox}
import pandas as pd
import numpy as np
import statsmodels.api as sm
import statsmodels.formula.api as smf

# 1. Load dataset
url = "https://stats.idre.ucla.edu/stat/stata/dae/nb_data.dta"
df = pd.read_stata(url)

# 2. Ensure categorical variable
df['prog'] = df['prog'].astype('category')

# 3. Fit Negative Binomial Model
model = smf.glm(
    'daysabs ~ C(prog) + math',
    data=df,
    family=sm.families.NegativeBinomial(alpha=1)
).fit()

print(model.summary())

# 4. IRR
irr = np.exp(model.params)
print("\n--- Incident Rate Ratios (IRR)---")
print(irr)

# 5. Overdispersion check
mean_days = df['daysabs'].mean()
var_days = df['daysabs'].var()

print(f"\nMean of daysabs: {mean_days:.4f}")
print(f"Variance of daysabs: {var_days:.4f}")

if var_days > mean_days:
    print("Overdispersion detected
    (Negative Binomial is appropriate)")
else:
    print("No strong overdispersion")
\end{codebox}

\textbf{Output:}
\begin{codebox}
--- Incident Rate Ratios (IRR)---
Intercept 13.667374
C(prog)[T.2.0] 0.643551
C(prog)[T.3.0] 0.278418
math 0.994030
Mean of daysabs: 5.9554
Variance of daysabs: 49.5187
Overdispersion detected (Negative Binomial is appropriate)
\end{codebox}

\textbf{Interpretation:}
The Negative Binomial regression model was applied due to overdispersion in the data, where the variance (49.52) is much greater than the mean (5.96).
This confirms that the Negative Binomial model is more appropriate than the Poisson model.
The coefficient for \textbf{math score} is negative and statistically significant ($p < 0.05$), indicating that higher math scores are associated with fewer days absent.
For program type, using Program 1 as the reference:
\begin{itemize}
\item Program 2 has a significant negative effect, meaning students have fewer absences compared to Program 1.
\item Program 3 has a stronger negative and highly significant effect, indicating substantially fewer absences.
\end{itemize}
Overall, both academic performance (math score) and program type significantly influence student absenteeism.

\begin{center}
\Large \textbf{Problem-8}
\end{center}
\textbf{Zero-Inflation Analysis of Count Data}
\textbf{Code:}
\begin{codebox}
import os
import pandas as pd
import numpy as np
import matplotlib.pyplot as plt
import seaborn as sns
import statsmodels.api as sm
import statsmodels.formula.api as smf

file_path = r"C:\Users\hp\Downloads\fish.csv"
save_dir = r"C:\Users\hp\OneDrive\Desktop\
Reality Stone-4.1\Siddik sir-4104"
sns.set_theme(style="whitegrid")

df_fish = pd.read_csv(file_path)

print("=== DATASET PREVIEW ===")
print(df_fish.head())

#. Zero percentage
zero_percentage = (df_fish['count'] == 0).mean() * 100
print(f"\nPercentage of zeros in 'count'
column: {zero_percentage:.2f}%")

# ৩ . Distribution plot
plt.figure(figsize=(7, 5))
sns.histplot(df_fish['count'], kde=False, 
color='teal', discrete=True)
plt.title(f"Distribution of Fish Caught\n(Excess Zeros: {zero_percentage:.2f}%)", fontsize=12, fontweight='bold')
plt.xlabel('Number of Fish Caught')
plt.ylabel('Frequency (Count)')
plt.savefig(os.path.join(save_dir,
'excess_zeros_distribution.png'), 
dpi=300, bbox_inches='tight')
plt.show()
\end{codebox}

\textbf{Output:}
\begin{codebox}
=== DATASET PREVIEW ===
nofish livebait camper persons child xb zg count
0 1 0 0 1 0 -0.896315 3.050405 0
1 0 1 1 1 0 -0.558345 1.746149 0
2 0 1 0 1 0 -0.401731 0.279939 0
3 0 1 1 2 1 -0.956298 -0.601526 0
4 0 1 0 1 0 0.436891 0.527709 1
Percentage of zeros in 'count' column: 56.80\%
\end{codebox}

\textbf{Graph:}
\begin{center}
\includegraphics[width=0.7\textwidth]{excess_zeros_distribution.png}
\end{center}

\textbf{Interpretation:}
The dataset shows that approximately 56.80\% of the observations have zero fish caught, indicating a high presence of excess zeros in the count variable.
The histogram further confirms that a large proportion of observations are concentrated at zero, with relatively fewer higher count values.
This suggests that a standard Poisson model may not be appropriate due to zero inflation.
Therefore, models such as \textbf{Zero-Inflated Poisson (ZIP)} or \textbf{Zero-Inflated Negative Binomial (ZINB)} would be more suitable for analyzing this type of data.

\textbf{Zero-Inflated Poisson (ZIP) Model}
\textbf{Code:}
\begin{codebox}
# Count model (Poisson part)
X_count = df_fish[['child', 'camper', 'persons']]
X_count = sm.add_constant(X_count)

# Inflation model (Logit part)
X_logit = df_fish[['child', 'persons']]
X_logit = sm.add_constant(X_logit)
y = df_fish['count']

# ২ . Zero-Inflated Poisson Model ফিট করা
zip_model = sm.ZeroInflatedPoisson(
    y, X_count,
    exog_infl=X_logit,
    inflation='logit'
).fit(maxiter=100)

print("\n=== ZERO-INFLATED POISSON (ZIP) MODEL SUMMARY ===")
print(zip_model.summary())
\end{codebox}

\textbf{Output:}
\begin{codebox}
=== ZERO-INFLATED POISSON (ZIP) MODEL SUMMARY ===
Pseudo R-squ.: 0.3296
inflate_const 0.9813
inflate_child 1.8227
inflate_persons -0.8454
const -0.8452
child -1.1472
camper 0.7597
persons 0.8337
\end{codebox}

\textbf{Interpretation:}
The Zero-Inflated Poisson (ZIP) model accounts for excess zeros in the dataset by combining a Poisson count model with a logit model for zero inflation.
The inflation (logit) part shows that:
\begin{itemize}
\item \textbf{child} has a positive and significant effect, indicating higher probability of structural zeros.
\item \textbf{persons} has a negative effect, meaning more people reduce the chance of zero outcomes.
\end{itemize}
In the count (Poisson) part:
\begin{itemize}
\item \textbf{child} has a significant negative effect on fish count.
\item \textbf{camper} has a positive and significant effect, increasing fish caught.
\item \textbf{persons} has a strong positive impact on fish count.
\end{itemize}
Overall, the model successfully captures both the excess zeros and the count process, making it more appropriate than a standard Poisson model.

\textbf{Incidence Rate Ratios (IRR) for Count Model}
\textbf{Code:}
\begin{codebox}
# Extract coefficients (first 4 parameters used as given)
count_params = zip_model.params[:4]

# Compute IRR
irr = np.exp(count_params)

# Create DataFrame
df_irr = pd.DataFrame({
    'Coefficients': count_params,
    'IRR (Rate Ratio)': irr
})

print("=== INCIDENCE RATE RATIOS (IRR) FOR COUNT MODEL ===")
print(df_irr)
\end{codebox}

\textbf{Output:}
\begin{codebox}
=== INCIDENCE RATE RATIOS (IRR) FOR COUNT MODEL ===
Coefficients IRR (Rate Ratio)
inflate_const 0.981276 2.667858
inflate_child 1.822725 6.188700
inflate_persons -0.845363 0.429402
const -0.845234 0.429457
\end{codebox}

\textbf{Interpretation:}
The Incidence Rate Ratios (IRR) represent the multiplicative effect of predictors on the expected count outcome.
An IRR greater than 1 indicates an increase in the expected count, while an IRR less than 1 indicates a decrease.
Here, \textbf{inflate\_child} has a strong positive effect (IRR ≈ 6.19), suggesting a large increase in the rate, whereas \textbf{inflate\_persons} and \textbf{const} have IRR values less than 1, indicating a reduction in the expected count.
Overall, the IRR transformation provides an easier interpretation of model coefficients in terms of rate changes.

\begin{center}
\Large \textbf{Problem-9}
\end{center}
\textbf{Exploratory Data Analysis}

\textbf{Code:}
\begin{codebox}
# 1. Load dataset
df_fish = pd.read_csv(file_path)

print("=== DATASET SHAPE & COLUMNS ===")
print(df_fish.shape)
print(df_fish.head())

# 2. Check percentage of zero values
zero_pct = (df_fish['count'] == 0).mean() * 100
print(f"\nPercentage of records with 0 fish caught: {zero_pct:.2f}%")

# 3. Plot distribution and save
plt.figure(figsize=(7, 5))
sns.histplot(df_fish['count'], color='dodgerblue', discrete=True, kde=False)
plt.title(f"Distribution of Fish Caught\n(Excess Zeros = {zero_pct:.2f}%)", fontsize=12, fontweight='bold')
plt.xlabel("Number of Fish Caught")
plt.ylabel("Frequency")
plt.savefig(os.path.join(save_dir, "fish_count_distribution.png"), dpi=300, bbox_inches='tight')
plt.show()
\end{codebox}

\textbf{Output:}
\begin{codebox}
=== DATASET SHAPE & COLUMNS ===
(250, 8)
Percentage of records with 0 fish caught: 56.80%
\end{codebox}

\textbf{Graph:}
\begin{center}
\includegraphics[width=0.7\textwidth]{123fish_count_distribution.png}
\end{center}

\textbf{Interpretation:}
The dataset contains 250 observations and 8 variables related to fishing activity.
A large proportion of the observations (56.80\%) have zero fish caught, indicating a high presence of excess zeros in the data.
The distribution plot further confirms that the data is highly skewed with a heavy concentration at zero.
This suggests that standard Poisson regression may not be appropriate, and a Zero-Inflated model (such as ZIP or ZINB) is more suitable for modeling this type of count data.

\textbf{Zero-Inflated Poisson (ZIP) Regression Analysis}
\textbf{Code:}
\begin{codebox}
# 1. Prepare variables for Count model
X_count = df_fish[['child', 'camper', 'persons']]
X_count = sm.add_constant(X_count)

# 2. Prepare variables for Zero-Inflation model
X_logit = df_fish[['child', 'persons']]
X_logit = sm.add_constant(X_logit)

# Target variable
y = df_fish['count']

# 3. Fit ZIP model
zip_model = sm.ZeroInflatedPoisson(
    y, X_count,
    exog_infl=X_logit,
    inflation='logit'
).fit(maxiter=100)

print("=== ZERO-INFLATED POISSON REGRESSION SUMMARY ===")
print(zip_model.summary())
\end{codebox}

\textbf{Output:}
\begin{codebox}
=== ZERO-INFLATED POISSON REGRESSION SUMMARY ===
Pseudo R-squ.: 0.3296
Log-Likelihood: -755.54
LLR p-value: 1.007e-160
Coefficients:
Inflation Model (Logit Part):
inflate_const = 0.9813 (p = 0.018)
inflate_child = 1.8227 (p < 0.001)
inflate_persons = -0.8454 (p < 0.001)
Count Model (Poisson Part):
const = -0.8452 (p < 0.001)
child = -1.1472 (p < 0.001)
camper = 0.7597 (p < 0.001)
persons = 0.8337 (p < 0.001)
\end{codebox}

\textbf{Interpretation:}
The Zero-Inflated Poisson (ZIP) model is used to handle excess zeros in the count data.
The model consists of two components:
\textbf{1. Inflation Model (Logit Part):} This part models the probability of always having zero fish caught.
The variable \textbf{child} has a strong positive effect, indicating that groups with children are more likely to catch zero fish.
The variable \textbf{persons} has a negative effect, suggesting that larger groups are less likely to belong to the zero-only category.
\textbf{2. Count Model (Poisson Part):} This part models the number of fish caught for groups that are actively fishing.
The variable \textbf{child} has a negative effect, meaning more children reduce the number of fish caught.
The variable \textbf{camper} has a positive effect, indicating campers tend to catch more fish.
The variable \textbf{persons} also has a positive effect, meaning larger groups catch more fish.
Overall, the model is highly significant (very small p-value), and the relatively high pseudo R-squared indicates a good fit.
The ZIP model is appropriate due to the presence of excess zeros in the dataset.

\textbf{Conversion of Coefficients to Incidence Rate Ratios (IRR)}
\textbf{Code:}
\begin{codebox}
# Extract coefficients from ZIP model
count_coefficients = zip_model.params[:4]

# Convert to IRR
irr_values = np.exp(count_coefficients)

# Create DataFrame
df_irr = pd.DataFrame({
    'Raw Coefficient': count_coefficients,
    'Incidence Rate Ratio (IRR)': irr_values
})

print("=== COEFFS CONVERTED TO INCIDENCE RATE RATIOS (IRR) ===")
print(df_irr)
\end{codebox}

\textbf{Output:}
\begin{codebox}
=== COEFFS CONVERTED TO INCIDENCE RATE RATIOS (IRR) ===
Raw Coefficient Incidence Rate Ratio (IRR)
inflate_const 0.981276 2.667858
inflate_child 1.822725 6.188700
inflate_persons -0.845363 0.429402
const -0.845234 0.429457
\end{codebox}

\textbf{Interpretation:}
The Incidence Rate Ratio (IRR) is obtained by exponentiating the model coefficients, allowing for easier interpretation in multiplicative terms. An IRR greater than 1 indicates an increase in the expected rate, while an IRR less than 1 indicates a decrease.
The variable \textbf{inflate\_child} shows a strong positive effect (IRR ≈ 6.19), meaning it substantially increases the likelihood associated with the inflation component.
On the other hand, \textbf{inflate\_persons} and \textbf{const} have IRR values less than 1, indicating a reduction in the expected rate.
Overall, transforming coefficients into IRR makes the model results more interpretable in terms of percentage change in rates.

\begin{center}
\Large \textbf{Problem-10}
\end{center}
\textbf{Negative Binomial Regression for Hospital Stay (Zero-Truncated Context)}
\textbf{Code:}
\begin{codebox}
# 1. Load dataset
df, meta = pyreadstat.read_dta("ztp.dta")
df.to_csv('ztp_temp.csv', index=False)
df = pd.read_csv('ztp_temp.csv')

# 2. Fit Negative Binomial GLM
model = smf.glm(
    'stay ~ age + hmo + died',
    data=df,
    family=sm.families.NegativeBinomial(alpha=1)
).fit()

# 3. Model summary
print(model.summary())

# 4. Compute IRR
irr = np.exp(model.params)
print("\n--- Incident Rate Ratios (IRR)---")
print(irr)
\end{codebox}

\textbf{Output:}
\begin{codebox}
Generalized Linear Model Regression Results
Model Family: NegativeBinomial
Pseudo R-squ. (CS): Approx. moderate fit
Coefficients:
Intercept = 2.4378 (p < 0.001)
age = -0.0147 (p = 0.371)
hmo = -0.1375 (p = 0.065)
died = -0.2038 (p < 0.001)
--- Incident Rate Ratios (IRR)---
Intercept = 11.4473
age = 0.9854
hmo = 0.8715
died = 0.8156
\end{codebox}

\textbf{Interpretation:}
A count regression model was applied to analyze hospital length of stay using a Negative Binomial specification, appropriate for overdispersed count data and consistent with a zero-truncated setting (no zero values in the response).
\textbf{Age:} The IRR is approximately 0.99 with a non-significant p-value ($p = 0.371$), indicating that age does not have a statistically significant effect on hospital stay duration.
\textbf{Health Insurance (HMO):} The IRR is 0.87, suggesting that patients with HMO insurance tend to have about 13\% shorter hospital stays.
However, the effect is only marginally significant ($p = 0.065$).
\textbf{Death Status (died):} The IRR is 0.82 and statistically significant ($p < 0.001$), indicating that patients who died had approximately 18\% shorter recorded hospital stays compared to those who survived.
\textbf{Conclusion:} Hospital stay duration is significantly influenced by death status, while age has no significant effect and HMO insurance shows a weak negative association.
The modeling approach is appropriate given the absence of zero values in the dependent variable.


\textbf{Zero-Truncated Negative Binomial (ZTNB) Regression Analysis}
\textbf{Code:}
\begin{codebox}
class ZTNB(GenericLikelihoodModel):
    def loglike(self, params):
        beta = params[:-1]
        alpha = params[-1]
        
        if alpha <= 0:
            return -np.inf
            
        mu = np.exp(np.dot(self.exog, beta))
        size = 1 / alpha
        prob = size / (size + mu)
        ll_nb = nbinom.logpmf(self.endog, size, prob)
        
        # Zero-truncation adjustment
        p0 = (1 + alpha * mu) ** (-1 / alpha)
        return np.sum(ll_nb - np.log(1 - p0))

df11 = pd.read_csv("ztp_temp.csv")
X11 = sm.add_constant(df11[['age', 'hmo', 'died']]).values
y11 = df11['stay'].values

start_p = np.append(
    sm.GLM(y11, X11, family=sm.families.Poisson()).fit().params,
    1.0
)

model11 = ZTNB(y11, X11).fit(start_params=start_p)
print(model11.summary())

print("\n--- IRRs ---")
print(np.exp(model11.params[:-1]))
\end{codebox}

\textbf{Output:}
\begin{codebox}
ZTNB Results
Log-Likelihood: -4755.3
AIC: 9521
No. Observations: 1493
Coefficients:
Intercept (const) = 2.4083 (p < 0.001)
age = -0.0157 (p = 0.231)
hmo = -0.1471 (p = 0.013)
died = -0.2177 (p < 0.001)
alpha = 0.5663 (p < 0.001)
--- IRRs ---
Intercept = 11.1150
age = 0.9844
hmo = 0.8633
died = 0.8043
\end{codebox}

\textbf{Interpretation:}
A Zero-Truncated Negative Binomial (ZTNB) regression model was used to analyze hospital length of stay, accounting for both overdispersion and the absence of zero values in the response variable.
\textbf{Age:} The IRR is approximately 0.98 with a non-significant p-value ($p = 0.231$), indicating that age does not significantly affect hospital stay duration.
\textbf{Health Insurance (HMO):} The IRR is 0.86 and statistically significant ($p = 0.013$), suggesting that patients with HMO insurance stay about 14\% fewer days compared to those without HMO.
\textbf{Death Status (died):} The IRR is 0.80 and highly significant ($p < 0.001$), indicating that patients who died had approximately 20\% shorter hospital stays than survivors.
\textbf{Dispersion Parameter ($\alpha$):} The estimated dispersion parameter is significant, confirming the presence of overdispersion and supporting the use of the Negative Binomial model over Poisson.
\textbf{Conclusion:} Hospital stay duration is significantly influenced by insurance type and death status, while age has no significant effect.
The ZTNB model is appropriate due to zero truncation and overdispersion in the data.

\end{document}

\begin{center}
\Large \textbf{Problem-11}
\end{center}